\begin{definition}[$\varepsilon$-сеть]\label{varepsilon-net}
    Пусть $M$~--- некоторое множество в метрическом пространстве $R$. Тогда множество $A$ из $R$ называется $\varepsilon$-сетью для $M$, если для любой точки $x \in M$ найдется хотя бы одна точка $a \in A$, что
    \[\rho(x, a) \leqslant \varepsilon.\]
\end{definition}

\begin{definition}[Ограниченость]
    Множество $M$ в метрическом пространстве $R$ называется ограниченым, если существует $B_\varepsilon(x_0)$ содержащий $M$.
\end{definition}

\begin{definition}[Вполне ограниченость, полная ограниченость]
    Множество $M$ в метрическом пространстве $R$ называется вполне ограниченым, если для него при любом $\varepsilon > 0$ существует \textbf{конечная} \nameref{varepsilon-net}.
\end{definition}

\begin{note}
    Если множество вполне ограничено, то его замыкагние также вполне ограничено.
\end{note}

\begin{note}
    Из определения следует, что если само пространство вполне ограничено, то оно сепарабельно.
\end{note}

\begin{note}
    Из вполне ограниченности следует ограниченость.
\end{note}
\begin{proof}
    Из вполне ограничености ограниченость получается как объединение конечного числа ограниченных множеств.
\end{proof}

\begin{note}
    В $n$-мерном евклидовом пространстве вполне ограниченость совпадает с обычной ограниченностью.
\end{note}

\begin{counterexample}[Из ограничености не следует вполне ограниченость]
    Рассмотрим единичную сферу $S$ в пространстве $l_2$. Выберем точки на $S$: 
    \begin{align*}
    e_1 &= (1, 0, 0, \ldots),\\
    e_2 &= (0, 1, 0, \ldots),\\
    e_3 &= (0, 0, 1, \ldots),\\
    \ldots
    \end{align*}
    Расстояние между любыми двумя отличными точками равно $\sqrt{2}$. Но тогда не существует конечной $\varepsilon$-сети при $\varepsilon < \frac{\sqrt{2}}{2}$.
\end{counterexample}

\begin{theorem}[Критерий компактности]\label{thm3.2}
    Пусть $X$ "--- метрическое пространство. Тогда следующие условия эквивалентны:
    \begin{enumerate}
        \item $X$ компактно
        \item $X$ полно и вполне ограниченно
        \item Из любой последовательности $\{x_n\} \subset X$ можно выделить сходящуюся подпоследовательность $\{x_{n_k}\}$ т.е. $X$ - \textit{секвенциально компактно}
        \item Любое бесконечное множество $M \subset X$ имеет предельную точку
    \end{enumerate}
\end{theorem}

\begin{proof}~
    \begin{itemize}
        \item\imp{1}{2}$X$ вполне ограниченно, поскольку для любого $\varepsilon > 0$ из открытого покрытия $\{B(x, \varepsilon)\}_{x \in X}$ по определению можно выделить конечное подпокрытие. Центры шаров из подпокрытия дадут требуемую $\varepsilon$-сеть. Остается доказать полноту пространства $X$.

        Пусть последовательность $\{x_n\} \subset X$ фундаментальна. Для каждого $n \in \N$ положим $A_n := \{x_n, x_{n+1}, \dotsc \}$, тогда система $\{\overline{A_n}\}$ является \hyperref[center-system]{центрированной системой замкнутых множеств}. Система центрирована потому что у любого конечного набора $\{\overline{A_n}\}_{n \in k_1, \ldots k_m}$ непустое пересечение: оно содержит хвост $\{x_{max\{k_1, \ldots k_m\}}, x_{max\{k_1, \ldots k_m\}+1}, \dotsc \}$.
        
        Поэтому можно выбрать точку $x_0 \in \bigcap\limits_{n \in \N_+} \overline{A_n},\; x_0 \in X$ по теореме~\ref{thm3.1}. В силу фундаментальности, для любого $\varepsilon > 0$ существует $N \in \N$ такое, что при всех $n > N$ выполнено $\overline{A_n} \subset \overline B(x_N, \varepsilon)$, откуда и $\rho(x_n, x_0) < \varepsilon$. Значит, $x_n \xrightarrow[X]{} x_0$.
        
        \item\imp{2}{3}Зафиксируем произвольную последовательность $\{x_n\} \subset X$. Поскольку $X$ вполне ограниченно, то для любого $\varepsilon > 0$ существует $y \in X$ такое, что шар $B(y, \varepsilon)$ содержит бесконечно много точек из $\{x_n\}$. Будем применять это рассуждение сначала для всего пространства $X$, потом для шаров в $X$, содержащих бесконечно много точек из $\{x_n\}$:
        \begin{itemize}
            \item[$\bullet$] Для $\varepsilon := 1$ выберем $\{x^1_k\} \subset \{x_n\}$ так, что $\{x^1_k\} \subset B(y_1, 1)$
            \item[$\bullet$] Для $\varepsilon := \frac 12$ выберем $\{x^2_k\} \subset \{x^1_n\} \subset B(y_1, 1)$ так, что $\{x^2_k\} \subset B(y_2, \frac 12)$
            \item[$\bullet$] Для $\varepsilon := \frac 13$ выберем $\{x^3_k\} \subset \{x^2_n\} \subset B(y_2, \frac 12)$ так, что $\{x^3_k\} \subset B(y_3, \frac 13)$, и так далее
        \end{itemize}

        Рассмотрим <<диагональную>> последовательность $\{x_k^k\} \subset \{x_n\}$. По построению, она является фундаментальной, и в силу полноты пространства $X$ она сходится.
        
        \item\imp{3}{1}Проверим сначала, что $X$ вполне ограниченно. Предположим противное, то есть существует $\varepsilon_0 > 0$ такое, что никакой конечный набор шаров радиуса $\varepsilon_0$ не покрывает $X$. Тогда можно выбрать точку $x_1 \in X$, затем точку $x_2 \in \big(X \bs B(x_1, \varepsilon_0)\big)$, точку $x_3 \in \big(X \bs (B(x_1, \varepsilon_0) \cup B(x_2, \varepsilon_0))\big)$, и так далее. Процесс не закончится, и будет получена последовательность $\{x_n\}$ с попарными расстояниями между точками не меньше $\varepsilon_0$, из которой нельзя выделить сходящуюся подпоследовательность, --- противоречие. Значит, $X$ вполне ограниченно.
        
        Теперь проверим, что $X$ компактно. Предположим противное, то есть существует открытое покрытие $\{G_\alpha\}_{\alpha \in \mf A}$, не имеющее конечного подпокрытия. Значит, для любого $\varepsilon > 0$ найдется $x \in X$ такое, что шар $B(x, \varepsilon)$ не покрывается конечным числом множеств из $\{G_\alpha\}$ (если такого шара нет, то из вполне ограниченности, складывая конечные покрытия конечного числа шаров, получим конечное покрытие всего множества). Выбирая такую точку $x_n$ для $\varepsilon := \frac 1n$ при каждом $n \in \N$, получим последовательность $\{x_n\}$, из которой можно выделить сходящуюся подпоследовательность $\{x_{n_k}\}$. Пусть $x_{n_k} \xrightarrow[X]{} x_0 \in X$, где $x_0 \in X$. Тогда существует $\alpha_0 \in \mf A$ такое, что $x_0 \in G_{\alpha_0}$. Но множество $G_{\alpha_0}$ открыто, поэтому оно покрывает некоторую окрестность точки $x_0$, а значит и все шары $B(x_{n_k}, \frac1{n_k})$, начиная с некоторого номера, --- противоречие.
        
        \item\imp{3}{4}Зафиксируем бесконечное множество $M \subset X$, тогда, выбирая произвольным образом последовательность $\{x_n\} \subset M$ и выделяя из нее сходящуюся подпоследовательность, получим требуемое.
        
        \item\imp{4}{3}Зафиксируем последовательность $\{x_n\}$. Если множество ее значений конечно, то в ней можно выделить стационарную подпоследовательность. Если же множество ее значений бесконечно, то оно имеет предельную точку $x_0 \in X$, поэтому можно выбрать подпоследовательность $\{x_{n_k}\}$ такую, что $x_{n_k} \xrightarrow[X]{} x_0$.\qedhere
    \end{itemize}
\end{proof}

\begin{note}
    Условию компактности пространства $X$ также эквивалентно такое условие: любая непрерывная функция $f : X \to \R$ является ограниченной на $X$, то есть выполнено включение $C(X, \R) \subset B(X, \R)$.
\end{note}

\begin{theorem}[Кантора]\label{thm3.3}
    Пусть $X$ "--- компактное МП, и функция $f: X \to \R$ непрерывна на $X$. Тогда $f$ равномерно непрерывна на $X$.
\end{theorem}

\begin{proof}
    Предположим противное, то есть выполнено следующее:
    \[\exists \varepsilon_0 > 0: \forall \delta > 0: \exists x, y \in X, \rho(x, y) < \delta: |f(x) - f(y)| \geqslant \varepsilon_0\]

    Выбирая $\delta := \frac 1n$ для каждого $n \in \N$, получим последовательности $\{x_n\}$ и $\{y_n\}$. Поскольку $X$ компактно, можно выделить из них сходящуюся подпоследовательности $\{x_{n_k}\}$ и $\{y_{n_k}\}$, причем сходятся они к одной и той же точке $x_0 \in X$ по построению, однако для любого $k \in \N$ выполнено $|f(x_{n_k}) - f(y_{n_k})| \geqslant \varepsilon_0$ --- противоречие.
\end{proof}

\begin{definition}
    Пусть $X$ "--- метрическое пространство. Множество $M \subset X$ называется \textit{предкомпактом}, если множество $\overline M$ компактно как подпространство в $X$.
\end{definition}

\begin{note}
    Пусть $M \subset X$. Нетрудно проверить, что справедливы следующие утверждения:
    \begin{itemize}
        \item Если множество $M$ вполне ограниченно как подмножество в $X$, то оно также вполне ограниченно как пространство, то есть конечную $\varepsilon$-сеть можно выбирать не в объемлющем пространстве $X$, а в самом пространстве $M$. Обратное, очевидно, тоже верно.
        
        \item Если замыкание $\overline M$ множества $M$ вполне ограниченно, то множество $M$ тоже вполне ограниченно.
        
        \item Если пространство $X$ полно, то условия полноты и замкнутости множества $M$ эквивалентны.
    \end{itemize}

    Из \hyperref[thm3.2]{критерия компактности} и утверждений выше следует, что в полном метрическом пространстве свойства предкомпактности и вполне ограниченности эквивалентны.
\end{note}